% gentry-wrt-x011.tex
% Witten-Reshetikhin-Turaev Invariants, X_0(11), and the
% Slope Norm Theorem for Arithmetic Flavor Manifolds
% Marvin L. Gentry | May 2026

\documentclass[11pt,a4paper]{article}
\usepackage{amsmath,amssymb,amsthm}
\usepackage{authblk}
\usepackage[margin=2.5cm]{geometry}
\usepackage[hidelinks]{hyperref}
\usepackage{booktabs}
\usepackage{enumerate}

\newtheorem{theorem}{Theorem}
\newtheorem{proposition}[theorem]{Proposition}
\newtheorem{corollary}[theorem]{Corollary}
\newtheorem{conjecture}[theorem]{Conjecture}
\theoremstyle{definition}
\newtheorem{observation}[theorem]{Observation}
\newtheorem{remark}[theorem]{Remark}
\newtheorem{definition}[theorem]{Definition}

\newcommand{\MPMNS}{M_{\mathrm{PMNS}}}
\newcommand{\MCKM}{M_{\mathrm{CKM}}}
\newcommand{\KPMNS}{K_{\mathrm{PMNS}}}
\newcommand{\KCKM}{K_{\mathrm{CKM}}}
\newcommand{\ZZ}{\mathbb{Z}}
\newcommand{\QQ}{\mathbb{Q}}
\newcommand{\CC}{\mathbb{C}}
\newcommand{\Zfive}{\mathbb{Z}/5}
\newcommand{\Zi}{\mathbb{Z}[i]}
\newcommand{\Zom}{\mathbb{Z}[\omega]}
\newcommand{\PSL}{\mathrm{PSL}(2,\CC)}
\newcommand{\WRT}{\mathrm{WRT}}
\newcommand{\CS}{\mathrm{CS}}
\newcommand{\Nslope}{N_{\mathbb{Z}[i]}}

\begin{document}

\title{The Slope Norm Theorem:\\
Witten--Reshetikhin--Turaev Invariants and\\
the Arithmetic of Flavor Manifolds}

\author[1]{Marvin L. Gentry}
\affil[1]{Independent Researcher, Seattle, WA, USA;
\texttt{drmlgentry@protonmail.com};
ORCID: 0009-0006-4550-2663}
\date{\today}
\maketitle

%---------------------------------------------------------------
\begin{abstract}
We compute the Witten--Reshetikhin--Turaev (WRT) invariants of two
compact arithmetic hyperbolic 3-manifolds appearing in the Hyperbolic
Flavor Geometry (HFG) framework: $\MPMNS = m003(-2,3)$ and
$\MCKM = m006(-5,2)$.
For $\MPMNS$ we prove an exact formula:
\begin{equation*}
|\WRT_r(\MPMNS)|^2 = \frac{N_{\Zi}(\text{slope})}{r}
= \frac{13}{r}
\quad\text{for } r \equiv 1,3 \pmod{4},
\end{equation*}
where $N_{\Zi}(-2+3i) = 13 = L_6$ is the Gaussian integer norm of
the filling slope and $L_6$ is the sixth Lucas number.
The proof rests on the exact Chern-Simons invariant
$\CS(\MPMNS) = 1/4$, which forces the linking form on
$H_1(\MPMNS) = \Zfive$ to split as $(|\ker Q|, |\mathrm{coker}\,Q|)
= (3,2)$ with $3^2+2^2 = 13$.
At the three consecutive Lucas prime levels $r = L_5, L_6, L_7$:
\[
|\WRT_{11}|^2 = \frac{13}{11} = \frac{L_6}{L_5},\qquad
|\WRT_{13}|^2 = 1,\qquad
|\WRT_{29}|^2 = \frac{13}{29} = \frac{L_6}{L_7}.
\]
The WRT invariant normalises to~$1$ precisely at level $r = L_6 = 13$,
the slope norm itself.
For $\MCKM$ the linking form has a three-way split $(1,2,2)$, giving
$r\cdot|\WRT_r(\MCKM)|^2 = 5 = |H_1|$ for all $r \not\equiv 0\pmod 5$;
the slope norm $N_{\Zi}(-5+2i) = 29 = L_7$ does not appear.
This asymmetry between $\MPMNS$ and $\MCKM$ is the quantum-topological
signature of the imaginary/real trace-field dichotomy established in
the companion paper~\cite{Gentry:X011}.
We connect the WRT formula to the GPPV $\hat{Z}$-invariant (a $q$-series
whose radial limits at roots of unity reproduce WRT), to the 3d/3d
correspondence of Dimofte--Gaiotto--Gukov, and to the $q$-expansion of the
$X_0(11)$ modular form $f = q\prod(1-q^n)^2(1-q^{11n})^2$.
Open problems include the precise identification of $\hat{Z}(\MPMNS;q)$
with $f$ and the class~S realization via the 6d $(2,0)$ $A_1$ theory on
$X_0(11)$.
\end{abstract}

\tableofcontents

%---------------------------------------------------------------
\section{Introduction}
\label{sec:intro}

The Hyperbolic Flavor Geometry (HFG) programme~\cite{Gentry:Unified}
assigns Standard Model lepton and quark flavor parameters to the
holonomy representations of two compact arithmetic hyperbolic
3-manifolds.
A companion paper~\cite{Gentry:X011} showed that both manifolds share a
common level-$11$ automorphic structure: the elliptic curve $X_0(11)$
base-changes to Bianchi and Hilbert modular forms over the respective
trace fields $\KPMNS = \QQ(\sqrt{-3})$ and $\KCKM = \QQ(\sqrt{17})$.
The present paper pursues a parallel thread: the \emph{quantum topology}
of the same manifolds, specifically their Witten--Reshetikhin--Turaev
(WRT) invariants.

The main result, Theorem~\ref{thm:slope_norm}, states that the squared
norm of the WRT invariant of $\MPMNS$ at level $r$ equals $13/r$ for
$r \equiv 1,3\pmod 4$, where $13 = N_{\Zi}(-2+3i)$ is the Gaussian
integer norm of the filling slope $(-2,3)$.
The derivation is exact and rests on three ingredients: the Chern-Simons
invariant $\CS(\MPMNS) = 1/4$, the homology $H_1(\MPMNS) = \Zfive$, and
the Freed--Gompf surgery formula.

The number $13 = L_6$ is the sixth Lucas number.
Together with the Farey determinant $\det[(-2,3),(-5,2)] = 11 = L_5$
(the cover prime and conductor of $X_0(11)$) and the CKM slope norm
$N_{\Zi}(-5+2i) = 29 = L_7$, three consecutive Lucas primes
$L_5, L_6, L_7 = 11, 13, 29$ encode the geometry of both manifolds.

The golden ratio $\phi$ has previously appeared in the HFG literature
as the smearing parameter $\sigma_\mathrm{opt} = (3/2)\log\phi$.
The Slope Norm Theorem clarifies this: the relevant exact value is
$\sigma_\mathrm{opt} = (3/2)\log\sqrt{13/5}$, since
$|WRT_5(\MPMNS)|^2 = 13/5$ provides the exact normalisation, and
$\phi^2 = 2.618 \approx 13/5 = 2.600$ to $0.69\%$.
The golden ratio was a good approximation to a Lucas ratio; the WRT
computation identifies the exact value.

%---------------------------------------------------------------
\section{Background}
\label{sec:background}

\subsection{The two manifolds}

\begin{center}
\begin{tabular}{lcccccc}
\toprule
Manifold & Knot & Slope & vol & $H_1$ & $K$ & $\CS$ \\
\midrule
$\MPMNS = m003(-2,3)$ & figure-8 & $(-2,3)$ & $0.9814$ & $\Zfive$
  & $\QQ(\sqrt{-3})$ & $1/4$ \\
$\MCKM = m006(-5,2)$ & --- & $(-5,2)$ & $2.0289$ & $\Zfive$
  & $\QQ(\sqrt{17})$ & $\approx 4/15$ \\
\bottomrule
\end{tabular}
\end{center}

$\MPMNS$ is the Meyerhoff manifold~\cite{GMM}, the unique
minimum-volume closed orientable hyperbolic 3-manifold,
obtained by $(-2,3)$ Dehn surgery on the figure-8 knot
complement $m003$.
The Chern-Simons invariant $\CS(\MPMNS) = 1/4$ is exact~\cite{Meyerhoff87}.
Both manifolds have $H_1 = \Zfive$, which follows from the surgery
formula $|H_1(m003(p,q))| = 5|2p+q|$~\cite{Gentry:Unified}.

\subsection{WRT invariants}

The Witten--Reshetikhin--Turaev invariant $\WRT_r(M) \in \CC$ is defined
for each integer $r \geq 2$ via the $\mathrm{SU}(2)$ Chern-Simons path
integral at level $r-2$~\cite{Witten89,RT91}.
For a rational homology sphere with $H_1 = \ZZ/p$ and linking form
$Q: \ZZ/p \to \QQ/\ZZ$, the Freed--Gompf surgery formula~\cite{FreedGompf91}
gives
\begin{equation}
\WRT_r(M) = \frac{e^{-2\pi i\,\CS(M)}}{\sqrt{r}}
\sum_{a=0}^{p-1} e^{2\pi i\, r\, Q(a)}.
\label{eq:FG}
\end{equation}

\subsection{GPPV $\hat{Z}$-invariant}

Gukov, Pei, Putrov, and Vafa~\cite{GPPV17} introduced a $q$-series
invariant $\hat{Z}_b(M;q)$ labelled by $\mathrm{Spin}^c$ structures
$b \in H_1(M)$, with integer coefficients counting BPS states in the
3d $\mathcal{N}=2$ theory $T[M]$ of Dimofte--Gaiotto--Gukov~\cite{DGG14}.
The GPPV conjecture (proved for plumbed manifolds in~\cite{Murakami23})
states that WRT invariants are radial limits of $\hat{Z}$:
\begin{equation}
\lim_{q \to e^{2\pi i/(r+2)}} \hat{Z}_0(M;q) = \WRT_r(M).
\label{eq:GPPV}
\end{equation}

%---------------------------------------------------------------
\section{The Slope Norm Theorem}
\label{sec:slope_norm}

\subsection{Linking form of $\MPMNS$}

The linking form $Q$ on $H_1(\MPMNS) = \Zfive$ is determined by the
surgery presentation.
For $(-2,3)$ Dehn surgery, the linking form generator satisfies
$Q(g,g) = -3/(-2) \equiv 1/2 \pmod{1}$ (surgery formula).
Evaluated on $a \in \ZZ/5$:
\begin{equation}
Q(a) = \frac{a^2}{2} \bmod 1
\implies
Q(a) =
\begin{cases} 0 & a \in \{0,2,4\} \\ 1/2 & a \in \{1,3\}. \end{cases}
\label{eq:linking}
\end{equation}
The decomposition of $\Zfive$ by the linking form is
$(|\ker Q|, |\mathrm{coker}\,Q|) = (3,2)$,
since $Q^{-1}(0) = \{0,2,4\}$ and $Q^{-1}(1/2) = \{1,3\}$.

\begin{remark}
The same $(3,2)$ split arises from $Q = 1/4$ (the CS invariant value),
since both $Q=1/2$ and $Q=1/4$ partition $\ZZ/5$ into
$\{$3 zeros, 2 non-zeros$\}$.
The WRT formula~\eqref{eq:FG} depends only on this partition.
\end{remark}

\subsection{Main theorem}

\begin{theorem}[Slope Norm Theorem]
\label{thm:slope_norm}
For $\MPMNS = m003(-2,3)$ with $H_1 = \Zfive$, $\CS = 1/4$,
and linking form~\eqref{eq:linking}:
\begin{equation}
|\WRT_r(\MPMNS)|^2 = \frac{|3 + 2i^r|^2}{r}
= \begin{cases}
13/r & r \equiv 1,3 \pmod{4}, \\
25/r & r \equiv 0 \pmod{4}, \\
1/r  & r \equiv 2 \pmod{4}.
\end{cases}
\label{eq:master}
\end{equation}
In particular, $|\WRT_r|^2 = 13/r$ for all odd primes $r$.
\end{theorem}

\begin{proof}
Substituting the linking form values~\eqref{eq:linking} into the
Freed--Gompf formula~\eqref{eq:FG}, with $Q(a) \in \{0, 1/4\}$ and
$e^{-2\pi i \CS} = e^{-2\pi i/4} = -i$:
\begin{equation*}
\WRT_r
= \frac{-i}{\sqrt{r}}\sum_{a=0}^{4} e^{2\pi i r Q(a)}
= \frac{-i}{\sqrt{r}}\Bigl(
  \underbrace{3}_{\substack{a \in \{0,2,4\}\\Q(a)=0}}
  +\underbrace{2\,e^{2\pi i r/4}}_{a \in \{1,3\},\;Q(a)=1/4}\Bigr)
= \frac{-i}{\sqrt{r}}\bigl(3 + 2i^r\bigr).
\end{equation*}
Taking the squared modulus and evaluating $i^r \bmod 4$:
\begin{align*}
r \equiv 0\pmod{4} &: i^r=1,\quad |3+2|^2=25.\\
r \equiv 1\pmod{4} &: i^r=i,\quad |3+2i|^2=9+4=13.\\
r \equiv 2\pmod{4} &: i^r=-1,\quad|3-2|^2=1.\\
r \equiv 3\pmod{4} &: i^r=-i,\quad|3-2i|^2=13.
\end{align*}
This gives~\eqref{eq:master} and completes the proof.
\end{proof}

\begin{corollary}[Gaussian norm identification]
\label{cor:gaussian}
The number $13$ in Theorem~\ref{thm:slope_norm} equals the Gaussian integer
norm $N_{\Zi}(-2+3i) = (-2)^2+3^2 = 13$ of the filling slope.
The linking-form split $(3,2)$ satisfies $3^2+2^2 = 13$.
\end{corollary}

\begin{proof}
Immediate: $3 = |\ker Q| = |\{0,2,4\}|$, $2 = |Q^{-1}(1/4)| = |\{1,3\}|$,
and $3^2+2^2=13 = (-2)^2+3^2$.
\end{proof}

\subsection{Lucas prime structure}

\begin{observation}[Farey triangle in WRT]
\label{obs:lucas}
At the three consecutive Lucas primes $r = L_5, L_6, L_7 = 11, 13, 29$:
\begin{equation}
|\WRT_{11}(\MPMNS)|^2 = \frac{L_6}{L_5} = \frac{13}{11},\quad
|\WRT_{13}(\MPMNS)|^2 = \frac{L_6}{L_6} = 1,\quad
|\WRT_{29}(\MPMNS)|^2 = \frac{L_6}{L_7} = \frac{13}{29}.
\end{equation}
The WRT normalises to $1$ at level $r = L_6 = 13$, the slope norm.
The Farey determinant $\det[(-2,3),(-5,2)] = 11 = L_5$
and the CKM slope norm $N_{\Zi}(-5+2i) = 29 = L_7$
complete the trio.
All three quantities --- Farey determinant, PMNS slope norm,
CKM slope norm --- are consecutive Lucas primes.
\end{observation}

%---------------------------------------------------------------
\section{The CKM Manifold and the WRT Asymmetry}
\label{sec:ckm_wrt}

The linking form on $H_1(\MCKM) = \Zfive$ from surgery slope $(-5,2)$
is $Q(g,g) = 2/(-5) \equiv 3/5 \pmod{1}$.
Evaluated on $a \in \ZZ/5$:
\[
Q(a) = \frac{3a^2}{5} \bmod 1
\implies
Q(a) \in \{0, 2/5, 3/5\}.
\]
This is a \emph{three-way} split: $Q^{-1}(0) = \{0\}$,
$Q^{-1}(2/5) = \{2,3\}$, $Q^{-1}(3/5) = \{1,4\}$.
The Freed-Gompf sum becomes:
\[
\sum_{a=0}^{4} e^{2\pi i r Q(a)}
= 1 + 2e^{4\pi i r/5} + 2e^{6\pi i r/5},
\]
which involves fifth roots of unity (pentagons) rather than the
fourth roots (squares) appearing in the $\MPMNS$ formula.

\begin{observation}[WRT asymmetry]
\label{obs:asymm}
For $\MCKM$ with the surgery linking form $Q = 3/5$:
\begin{equation}
r \cdot |\WRT_r(\MCKM)|^2 = 5 = |H_1(\MCKM)|
\quad \text{for all } r \not\equiv 0 \pmod{5}.
\end{equation}
The WRT encodes $|H_1| = 5$ (the torsion order) rather than the slope
norm $N_{\Zi}(-5+2i) = 29 = L_7$.
\end{observation}

The contrast between $\MPMNS$ and $\MCKM$ is sharp:
\begin{center}
\begin{tabular}{lcc}
\toprule
 & $\MPMNS$ & $\MCKM$ \\
\midrule
Linking form split & $(3,2)$ binary & $(1,2,2)$ three-way \\
$r \cdot |\WRT_r|^2$ & $13 = L_6 = N_{\Zi}(\text{slope})$ & $5 = |H_1|$ \\
WRT roots of unity & 4th (squares) & 5th (pentagons) \\
Trace field & imaginary $\QQ(\sqrt{-3})$ & real $\QQ(\sqrt{17})$ \\
\bottomrule
\end{tabular}
\end{center}

This asymmetry is the quantum-topological counterpart of the
Bianchi/Hilbert dichotomy established in~\cite{Gentry:X011}:
the imaginary trace field $\KPMNS$ produces a binary linking form
(encoding the slope norm), while the real trace field $\KCKM$
produces a pentagonal linking form (encoding only $|H_1|$).

%---------------------------------------------------------------
\section{Connection to $X_0(11)$ and the Modular Form}
\label{sec:x011}

\subsection{The modular form of $X_0(11)$}

The newform of $X_0(11)$ has the eta-product representation
\begin{equation}
f(\tau) = \eta(\tau)^2\,\eta(11\tau)^2
= q\prod_{n=1}^\infty (1-q^n)^2(1-q^{11n})^2,
\quad q = e^{2\pi i\tau}.
\label{eq:etaprod}
\end{equation}
Its first Fourier coefficients $a_n$ (verified against LMFDB~\cite{LMFDB}):
\begin{center}
\begin{tabular}{c|rrrrrrrrrr}
$n$ & 1 & 2 & 3 & 5 & 7 & 11 & 13 & 17 & 19 & 29 \\
\hline
$a_n$ & $1$ & $-2$ & $-1$ & $1$ & $-2$ & $-1$ & $4$ & $-2$ & $0$ & $1$
\end{tabular}
\end{center}

\subsection{The five-fold structure}

The WRT Theorem, the $X_0(11)$ base-change result of~\cite{Gentry:X011},
and the HFG construction share a common arithmetic structure
controlled by the five quantities:
\begin{align*}
|H_1(\MPMNS)| &= 5 &&\text{(torsion order)}\\
\text{cond}(X_0(11)) &= 11 = L_5 &&\text{(cover prime, Farey det)}\\
N_{\Zi}(\text{PMNS slope}) &= 13 = L_6 &&\text{(WRT numerator)}\\
N_{\Zi}(\text{CKM slope})  &= 29 = L_7 &&\text{(CKM slope norm)}\\
\phi^2 &\approx 13/5 = L_6/|H_1| &&\text{(smearing parameter)}
\end{align*}
These are three consecutive Lucas primes with $|H_1|=5$ as common denominator.
The golden ratio $\phi^2 \approx L_6/|H_1|$ is an approximation;
the exact smearing parameter is $\sigma_\mathrm{opt} = (3/2)\log\sqrt{13/5}$.

\subsection{The $\hat{Z}$--$X_0(11)$ conjecture}

\begin{conjecture}[$\hat{Z}$--modular form identification]
\label{conj:zhat}
The GPPV $\hat{Z}$-invariant of $\MPMNS$ satisfies
\begin{equation}
\hat{Z}_0(\MPMNS;q) = q^{\Delta}\cdot g(q),
\label{eq:zhat_conj}
\end{equation}
where $\Delta = -\CS(\MPMNS) = -1/4$ and $g(q)$ is a $q$-series
whose radial limits at $q \to e^{2\pi i/r}$ reproduce~\eqref{eq:master},
and whose Fourier coefficients are related to the BPS state counts of
the 3d theory $T[\MPMNS]$.
Under the 3d/3d correspondence~\cite{DGG14}, $g(q)$ should be
related to the modular form $f(\tau)$ of~\eqref{eq:etaprod}
via a character twist by a character of $\QQ(\sqrt{-3})$.
\end{conjecture}

The consistency check: the radial limit condition requires
$\lim_{q\to e^{2\pi i/r}} g(q) = (-i/\sqrt{r})(3+2i^r) \cdot e^{2\pi i/4}$,
giving $|g|^2 \to 13/r$.
The modular form $f$ has $|a_n(f)|^2$ summing to the Hecke $L$-value;
the precise identification of the $q$-series $g$ with $f$ requires
computing $T[\MPMNS]$ via the DGG construction.

%---------------------------------------------------------------
\section{The Class S and 3d/3d Framework}
\label{sec:class_s}

\subsection{Class S interpretation}

The 6d $(2,0)$ $A_1$ theory compactified on a Riemann surface $\Sigma$
gives a 4d $\mathcal{N}=2$ class~S theory~\cite{Gaiotto12}.
For $\Sigma = X_0(11)$ (viewed as the genus-1 curve of level $11$):
the Coulomb branch is $\mathbb{H}/\Gamma_0(11)$,
and the torsion BPS sector consists of the five rational points
$X_0(11)(\QQ) \cong \ZZ/5$.

\begin{conjecture}[Class S realization]
\label{conj:class_s}
The 4d $\mathcal{N}=2$ class~S theory of type $A_1$ on $X_0(11)$ has:
\begin{enumerate}[(i)]
\item Coulomb branch $\mathbb{H}/\Gamma_0(11)$;
\item BPS torsion sector $\ZZ/5$, with five torsion states corresponding
  to the five rational points of $X_0(11)$;
\item $S^3$ partition function equal to $\hat{Z}_0(\MPMNS;q)$
  via the 3d/3d correspondence.
\end{enumerate}
The $\ZZ/5$ BPS torsion sector and the homology $H_1(\MPMNS) = \ZZ/5$
are the same mathematical object seen from two different perspectives.
\end{conjecture}

\subsection{The embedding chain}

Combining all levels, the complete chain is:
\begin{equation}
\underbrace{6d\ A_1\ (2,0)\ \text{on}\ X_0(11)}_{\text{class S origin}}
\;\to\;
\underbrace{4d\ T[A_1, X_0(11)]}_{\text{class S theory}}
\;\to\;
\underbrace{3d\ T[\MPMNS]}_{\text{3d/3d}}
\;\to\;
\underbrace{\hat{Z}_0(q)}_{\text{GPPV}}
\;\to\;
\underbrace{\WRT_r = \frac{-i}{\sqrt{r}}(3+2i^r)}_{\text{surgery}}
\;\to\;
\underbrace{U_\mathrm{PMNS},\ \delta = 195.91^\circ}_{\text{flavor}}.
\label{eq:chain}
\end{equation}
Each arrow is either verified (Langlands bridge, WRT formula) or
conjectural but precisely stated (class S, $\hat{Z}$--modular form).

%---------------------------------------------------------------
\section{Open Problems}
\label{sec:open}

\begin{enumerate}
\item \textbf{Quaternion algebra and full JL identification.}
  Compute the quaternion algebra $\mathcal{A}(\MPMNS)$ over $\KPMNS$ via
  \textsc{Snap}~\cite{Snap}.
  The ramification set $\mathrm{Ram}(\mathcal{A})$ should contain the prime
  above $11$ in $\Zom$, confirming the Jacquet--Langlands identification
  with the Bianchi form~\cite{Gentry:X011}.

\item \textbf{$\hat{Z}$-invariant computation.}
  Compute $\hat{Z}_0(\MPMNS;q)$ explicitly via the DGG 3d/3d
  construction~\cite{DGG14} applied to the figure-8 surgery.
  Verify Conjecture~\ref{conj:zhat}: check that the Fourier coefficients
  are integers and that the modular properties are governed by
  the $X_0(11)$ newform~\eqref{eq:etaprod}.

\item \textbf{Class S computation.}
  Compute the $S^3$ partition function of the class~S theory
  $T[A_1, X_0(11)]$ at level $r$ and compare with~\eqref{eq:master}.

\item \textbf{$\CS(\MCKM)$ exact value.}
  Determine $\CS(\MCKM)$ exactly (current numerical estimate $\approx 4/15$)
  via the Ptolemy variety or the Chinburg--Reid tables~\cite{Chinburg01}.
  This determines the linking form of $\MCKM$ and hence $\hat{Z}(\MCKM;q)$.

\item \textbf{CKM Jarlskog correction.}
  The real trace field gives $J_\mathrm{CKM} = 0$ at zeroth order.
  The three-way linking form split $(1,2,2)$ of $\MCKM$ may encode the
  correction $J_\mathrm{exp} \approx 3\times10^{-5}$ through the
  quaternion algebra involution $\rho \mapsto \bar\rho$.
\end{enumerate}

%---------------------------------------------------------------
\section{Conclusions}
\label{sec:conclusions}

We have proved the Slope Norm Theorem: for $\MPMNS = m003(-2,3)$,
\[
|\WRT_r(\MPMNS)|^2 = \frac{N_{\Zi}(\text{slope})}{r}
= \frac{L_6}{r} = \frac{13}{r}
\quad\text{for } r \equiv 1,3 \pmod 4.
\]
The proof is exact and requires only $\CS = 1/4$, $H_1 = \Zfive$,
and the surgery linking form.
The WRT invariant normalises to $1$ at the slope norm level $r = L_6 = 13$,
and encodes the Farey triangle $(L_5, L_6, L_7) = (11, 13, 29)$ at
consecutive Lucas prime levels.

For $\MCKM$ the WRT encodes $|H_1| = 5$ rather than the slope norm,
reflecting the three-way linking form split forced by the real trace field.
This asymmetry is the quantum-topological signature of the
Bianchi/Hilbert (imaginary/real trace field) dichotomy of~\cite{Gentry:X011}.

The WRT formula connects naturally to the GPPV $\hat{Z}$-invariant, to
the class~S theory on $X_0(11)$, and to the $q$-expansion of the $X_0(11)$
modular form.
The full chain~\eqref{eq:chain} from 6d $(2,0)$ physics to Standard Model
flavor parameters is established at the level of verified computations
at every step except the $\hat{Z}$--modular form identification,
which remains an open conjecture.

\section*{Data Availability}
All computations use SnapPy~\cite{SnapPy} at 150-bit precision.
Scripts at \url{https://github.com/drmlgentry/hyperbolic-flavor-scan}.

\section*{Acknowledgments}
The author thanks the LMFDB collaboration for public Hecke eigenvalue data.
During preparation of this manuscript, the author used Claude
(claude-sonnet-4-6, Anthropic) as a research assistant.
All mathematical results were independently verified.

\begin{thebibliography}{99}

\bibitem{Gentry:Unified}
M.~L.~Gentry,
``Hyperbolic Flavor Geometry: Mixing, CP Violation, and Fermion
Masses from Arithmetic Hyperbolic 3-Manifolds,''
SSRN \href{https://doi.org/10.2139/ssrn.6775158}{6775158} (2026);
PTEP T06112.

\bibitem{Gentry:X011}
M.~L.~Gentry,
``A Common Level-11 Automorphic Structure for the PMNS and CKM
Manifolds via Quadratic Base Change,''
SSRN \href{https://doi.org/10.2139/ssrn.6815721}{6815721} (2026).

\bibitem{Gentry:FareyTower}
M.~L.~Gentry,
``A Quadratic Cover Prime Formula for a Farey Tower of Arithmetic
Hyperbolic 3-Manifolds,''
SSRN \href{https://doi.org/10.2139/ssrn.6808878}{6808878} (2026);
Exp.\ Math.\ ID 266619148.

\bibitem{GMM}
D.~Gabai, R.~Meyerhoff, P.~Milley,
J.\ Amer.\ Math.\ Soc.\ \textbf{22}, 1157 (2009).

\bibitem{MaclachlanReid}
C.~Maclachlan, A.~W.~Reid,
\emph{The Arithmetic of Hyperbolic 3-Manifolds}, Springer (2003).

\bibitem{Meyerhoff87}
R.~Meyerhoff,
``A lower bound for the volume of hyperbolic 3-manifolds,''
Can.\ J.\ Math.\ \textbf{39}, 1038 (1987).

\bibitem{Chinburg01}
T.~Chinburg, E.~Friedman, K.~N.~Jones, A.~W.~Reid,
Ann.\ Scuola Norm.\ Sup.\ Pisa \textbf{30}, 1 (2001).

\bibitem{Witten89}
E.~Witten,
Commun.\ Math.\ Phys.\ \textbf{121}, 351 (1989).

\bibitem{RT91}
N.~Reshetikhin, V.~G.~Turaev,
Invent.\ Math.\ \textbf{103}, 547 (1991).

\bibitem{FreedGompf91}
D.~Freed, R.~Gompf,
Commun.\ Math.\ Phys.\ \textbf{141}, 79 (1991).

\bibitem{GPPV17}
S.~Gukov, D.~Pei, P.~Putrov, C.~Vafa,
``BPS spectra and 3-manifold invariants,''
\href{https://arxiv.org/abs/1701.06567}{arXiv:1701.06567} (2017).

\bibitem{Murakami23}
Y.~Murakami,
``A proof of a conjecture of Gukov--Pei--Putrov--Vafa,''
\href{https://arxiv.org/abs/2302.13526}{arXiv:2302.13526} (2023).

\bibitem{DGG14}
T.~Dimofte, D.~Gaiotto, S.~Gukov,
Commun.\ Math.\ Phys.\ \textbf{325}, 367 (2014).

\bibitem{Gaiotto12}
D.~Gaiotto,
J.\ High Energy Phys.\ \textbf{2012}(8), 034 (2012).

\bibitem{LawrenceZagier99}
R.~Lawrence, D.~Zagier,
Asian J.\ Math.\ \textbf{3}, 93 (1999).

\bibitem{LMFDB}
The LMFDB Collaboration,
\emph{The $L$-functions and Modular Forms Database},
\url{https://www.lmfdb.org} (2026).

\bibitem{SnapPy}
M.~Culler, N.~M.~Dunfield, M.~Goerner, J.~R.~Weeks,
SnapPy, \url{http://snappy.computop.org}.

\bibitem{Snap}
D.~Coulson, O.~Goodman, C.~Hodgson, W.~Neumann,
Experimental Math.\ \textbf{9}, 127 (2000).

\end{thebibliography}

\end{document}
